% !TeX root = ../papers/nc-p3.tex
\section{孤立分岐点の二側切断と後続の有理正規曲線}
\label{sec:isolated-branch-blocks}
\begin{proposition}\label{prop:isolated-branch-blocks}
位置0の点 $p$ で順方向支持 $S\subset V_1$ が二次元とする。逆支持を $T\subset V_3$、$t=\dim T$ とする。
$t=2,3$ に限られ、$R_1\cap(S\otimes V_2)$ は $w(p)$ の末尾縮約像に等しく次元 $t$ である。
後続直線 $L=\mathbb P(S^\perp)$ の一般点からさらに点が続く場合、それは一定点、
または $L$ と射影同型にパラメータ付けられる直線、滑らかな二次曲線、有理正規三次曲線である。
$t=3$ なら一般点に後続があり、$L$ 上で順方向階数が落ちるスキーム
（三次小行列式の共通零点）の長さは $3-d$ である。ここで $d$ はこのパラメータ表示の次数である。
$t=2$ で一般点が後続を持たないときは、後続を持つ点が二項四次の零因子をなす。
\end{proposition}
\begin{proof}
$t=4$ なら辺 $(p,q)$ の直交不等式により全 $q\in L$ で順方向支持が二次元となり
定理\ref{thm:no-moving-plane-pencils}に反する。
交叉次元は $t$ 以上である。一方 $L$ 上で残る二次式は一般点で階数3以上なので、交叉次元は3以下である。
$t=2$ のときは系\ref{cor:endpoint-line-line-cut}でちょうど2である。

一般に階数3のときは、$L=\mathbb P^1$ 上の小行列式が定める核点への有理写像を延長し、
像の文字の張る空間を $W\subset H^0(\mathcal O_L(d))$ とする。共通因子を除いてあるので基点はない。
$W$ の評価核を $N$ とする。位置2の全四文字から $W$ に落としても、
二次関係は $N(1)$ の一般ファイバーを張る。もとの点評価が後続点の全消滅超平面だからである。
$N$ の分裂に自明または正次数の因子はない（全切断の線型独立性による）。
また一般点での生成性により $-2$ 以下の因子もない。従って
$N=\mathcal O(-1)^{\dim W-1}$ で $d=\dim W-1$、$W$ は完全系である。
$t=3$ では $L$ に制限した三式の評価が一般に階数3であり、核座標は三次小行列式だけで書ける。
その共通因子の次数は $3-d$ である。$t=2$ で一般階数4の場合、制限した四式の行列式は非零な四次である。
\end{proof}

\begin{proposition}\label{prop:two-sided-center-self-block}
$t=2$ の点では $d(p,p)=1$（定理\ref{thm:recurrent-point-correspondences}の三文字商）であり、
$X=V_1/S$, $F=V_3/T$ に制限した四テンソルは $a\otimes h$、
$0\ne a\in\ker p$, $0\ne h\in X\otimes V_2\otimes F$ となる。
前後どちらの隣接直線にも、その辺で逆向きのファイバーが直線になる点が存在する。
より正確には、$H=a^\perp\subset V_0^*$ とすると、順側で $H$ 内に $p$ 以外の前点を持つ軌跡は
非零二項二次式の零因子の台であり、逆側でも同様である。
\end{proposition}
\begin{proof}
制限した $w$ の位置0対他の階数は高々1である。零なら任意の $u\in F^*$ について
$w(u)$ は位置1で $S$ に支持される。従って $V_0$ からの縮約像は
$R_1\cap(S\otimes V_2)$ に入り、高々二次元である。すると全 $\mathbb P(F^*)$ で
位置0の端支持が二次元となり可動中心が生じる。よって階数は1で、表示の $a$ は $w(p)$ の制限が零なので $\ker p$ に入る。
位置0からのこの非零縮約は二つの内側関係の挿入の交叉に入る。
両挿入は各八次元、周囲は十六次元なので交叉と余核は同じ次元であり、
反復ブロックの多項式増大で1以下となる。従って両者とも一次元である。

$q\in\mathbb P(X^*)$ について $H$ を位置0、$q$ を位置1で縮約すれば $F$ 側への像が零となる。
従って値は $R_2\cap(V_2\otimes T)$ の二空間に入り、常に $p$ が核にある。
これは $H/kp$ からこの二空間への $q$ に線型な二次正方行列である。
一般の $q$ では前点は $p$ 一つなので一般に可逆であり、行列式は非零二次である。
その根では $H$ 内の前点が少なくとも直線をなし、射影平面ファイバーはないので直線になる。
語順を反転すれば逆側での主張を得る。
\end{proof}
